% The reserved notation, carried on the tetrahedron at (4,3) with the triple
% C = {0,1,2}.  Four panels: the design in R^3 with its per-atom scalars; the
% two products of M_C on R^3; the chart pipeline V -> P -> W on the index space
% R^4, with the block W[C] marked; and a size strip placing the design against
% the window N(3).
%
% Orthographic projection of R^3 with azimuth 15 and elevation 70 degrees,
% scaled by 1.10cm:
%   X = 1.10(-0.2588190 x_1 + 0.9659258 x_2)
%   Y = 1.10(-0.9076734 x_1 - 0.2432103 x_2 + 0.3420201 x_3)
% This is the projection of figures/tetrahedron.tex and figures/tetragap.tex,
% unaltered in direction and reduced from 1.35cm to 1.10cm so that the panels
% fit the measure; the three figures are therefore registered up to the scale
% factor 1.10/1.35.  The viewer lies along e = X x Y =
% (0.3303661, 0.0885213, 0.9396926), a larger e-component being nearer; the four
% atoms sit at e-components g_0 = +1.3586, g_3 = +0.5208, g_1 = -0.6978,
% g_2 = -1.1815.  The cube is copied from figures/tetrahedron.tex and its twelve
% edges are drawn alike, with no hidden-line distinction, exactly as there; the
% one hidden-line decision made here is that the shaded triangle on the tips of
% g_0, g_1, g_2 is laid down before the arrows, so the arrows cross it.
%
% EXACT DATA, from Example ex:tetra and Appendix app:tetra.  Atoms and weights
%   g_0 = ( 1, 1, 1)   g_1 = ( 1,-1,-1)   g_2 = (-1, 1,-1)   g_3 = (-1,-1, 1)
%   t_c = 1/4 for every c,  sum_c t_c = 1.
% Each diagonal entry of sum_c g_c g_c^T is 4 and each off-diagonal entry is
% 1-1-1+1 = 0, so sum_c g_c g_c^T = 4 I_3 and sum_c t_c g_c g_c^T = I_3 over the
% rationals: a design at (4,3).  Hence, all over the rationals,
%   ell_c = 3,  s_c = t_c ell_c = 3/4,  sum_c s_c = 3 = k,
%   v_c = sqrt(t_c) g_c = g_c/2,  |v_c|^2 = s_c = 3/4,
%   V = (1/2) M with M the 4x3 matrix of rows g_c^T;  V^T V = I_3,
%   M M^T = 4 I_4 - J_4, so P = V V^T = (1/4)(4 I_4 - J_4) = I_4 - J_4/4,
%   P_cc = 3/4 = s_c,  P_cd = -1/4 = (1/4)<g_c,g_d>,  P^2 = P,  tr P = 3 = k,
%   D = (1/4) I_4,  W = P - D = (1/4)(3 I_4 - J_4),
%   W_cc = 1/2 = t_c(ell_c - 1),  sum_c W_cc = 2 = k - 1.
% At C = {0,1,2}, with M_C the 3x3 matrix of rows g_0^T, g_1^T, g_2^T,
%   det M_C = 4,
%   S_C    = M_C^T M_C = 4 I_3 - g_3 g_3^T = [[3,-1,1],[-1,3,1],[1,1,3]],
%   Gram_C = M_C M_C^T = 4 I_3 - J_3       = [[3,-1,-1],[-1,3,-1],[-1,-1,3]],
%   both of trace 9 and determinant 16 and both of spectrum {1,4,4},
%   W[C] = (1/4)(3 I_3 - J_3),  spectrum {0, 3/4, 3/4},  det W[C] = 0,
%   P[C] = (1/4)(4 I_3 - J_3),  pi_C = det P[C] = 16/64 = 1/4.
% Two cross-checks are printed on the plate.  The determinant of W[C] is taken
% twice: directly it is 0, and through cor:minor it is
% (prod_{c in C} t_c) det(Gram_C - I_3) = (1/64) det(3 I_3 - J_3) = (1/64)*0;
% the second reading applied to the pair {0,1} gives
% det W[{0,1}] = (1/16) det [[2,-1],[-1,2]] = 3/16, so the leading 2x2 minor is
% positive and only the 3x3 one vanishes.  And the two products of M_C are
% compared: different matrices, equal trace 9, equal determinant 16, equal
% spectrum {1,4,4} -- which is lem:flip at the one selection size, |C| = k, at
% which M_C is square.
% The window at rank three is N(3) = 3*4/2 + 1 = 7, and the design sits at m = 4.
%
% DISCRETIONARY CHOICES.  (i) The objects on the index space R^4 -- V, P, W --
% are drawn as cell arrays because the point of that panel is a principal block,
% and only a cell array can carry one; M_C, S_C and Gram_C are drawn as
% matrices.  The two products do not act on one space: S_C = M_C^T M_C is the
% k x k statement on the ambient R^3, while Gram_C = M_C M_C^T is the |C| x |C|
% statement on the block C of the index space, and the panel heading names the
% two sizes rather than one space (rem:size).  (ii) The block W[C] is filled black!25 rather than the
% black!38 of figures/chart.tex because here it carries numerals.  (iii) Only
% column indices are printed on P and W: both are symmetric, so the row indices
% repeat them.  (iv) Indices run 0..3, as in Example ex:tetra, not 1..m.
% (v) Every array on the index space is an integer array times 1/2 or 1/4, and
% is printed that way, so that no entry is a compound fraction.  (vi) The design
% is real, so every adjoint here is a transpose and the plate writes T where the
% table of Section 1 writes H.
%
% DISCLOSED.  The tetrahedron is uniform in all three per-atom scalars: t_c,
% ell_c and s_c are constant in c, so D is a scalar matrix here and the plate
% cannot show a leverage score varying across the atoms, which is the general
% case and the case Section 7 works in.  It is also a design at which every
% triple sits on the threshold, so the plate shows domination with slack zero
% and never with slack positive.
\begin{tikzpicture}[
  x=1cm, y=1cm,
  heavy/.style={-{Latex[length=2mm,width=1.5mm]}, line width=1.05pt},
  vecomit/.style={-{Latex[length=2mm,width=1.5mm,fill=white]},
                  line width=0.75pt, black!55},
  cube/.style={densely dotted, line width=0.3pt, black!35},
  ar/.style={-{Latex[length=2mm,width=1.5mm]}, line width=0.5pt},
  cellline/.style={draw=black!45, line width=0.25pt},
  border/.style={draw=black!85, line width=0.5pt},
  every node/.style={font=\small}]

% a grid of #3 columns and #4 rows, upper left corner (#1,#2), cell side 0.44
\newcommand{\cellframe}[4]{%
  \foreach \gi in {0,...,#3}{%
    \draw[cellline] ({#1+0.44*\gi},#2) -- ({#1+0.44*\gi},{#2-0.44*#4});}
  \foreach \gj in {0,...,#4}{%
    \draw[cellline] (#1,{#2-0.44*\gj}) -- ({#1+0.44*#3},{#2-0.44*\gj});}
  \draw[border] (#1,#2) rectangle ({#1+0.44*#3},{#2-0.44*#4});}
% the entries, row by row, upper left corner (#1,#2)
\newcommand{\cellnums}[3]{%
  \foreach \rowlist [count=\ri from 0] in {#3}{%
    \foreach \val [count=\ci from 0] in \rowlist {%
      \node[font=\scriptsize, inner sep=0pt]
        at ({#1+0.44*\ci+0.22},{#2-0.44*\ri-0.22}) {$\val$};}}}
% column indices 0..3 over a grid with upper left corner (#1,#2)
\newcommand{\cellindex}[2]{%
  \foreach \ii in {0,1,2,3}{%
    \node[font=\tiny, inner sep=1pt] at ({#1+0.44*\ii+0.22},{#2+0.14}) {\ii};}}

% ==================== panel A: the design in R^3 =========================
\node[anchor=south, font=\small] at (1.85,4.73) {$\dsn$ at $m=4$, $k=3$};
\begin{scope}[shift={(1.85,3.05)}]
  % the eight vertices of [-1,1]^3, at scale 1.10
  \coordinate (ppp) at ( 0.7778,-0.8897);   % g_0 = ( 1, 1, 1)
  \coordinate (ppm) at ( 0.7778,-1.6422);   %       ( 1, 1,-1)
  \coordinate (pmp) at (-1.3472,-0.3547);   %       ( 1,-1, 1)
  \coordinate (pmm) at (-1.3472,-1.1071);   % g_1 = ( 1,-1,-1)
  \coordinate (mpp) at ( 1.3472, 1.1071);   %       (-1, 1, 1)
  \coordinate (mpm) at ( 1.3472, 0.3547);   % g_2 = (-1, 1,-1)
  \coordinate (mmp) at (-0.7778, 1.6422);   % g_3 = (-1,-1, 1)
  \coordinate (mmm) at (-0.7778, 0.8897);   %       (-1,-1,-1)
  \coordinate (O)   at ( 0,0);

  \draw[cube] (ppp)--(ppm) (ppp)--(pmp) (ppp)--(mpp)
              (ppm)--(pmm) (ppm)--(mpm) (pmp)--(pmm) (pmp)--(mmp)
              (mpp)--(mpm) (mpp)--(mmp) (pmm)--(mmm) (mpm)--(mmm) (mmp)--(mmm);

  % the face on the tips of the selected triple
  \path[fill=black!12] (ppp)--(pmm)--(mpm)--cycle;

  \draw[heavy]   (O)--(ppp) node[below right=-1pt] {$g_0$};
  \draw[heavy]   (O)--(pmm) node[below left=-1pt]  {$g_1$};
  \draw[heavy]   (O)--(mpm) node[right=1pt]        {$g_2$};
  \draw[vecomit] (O)--(mmp)
    node[anchor=east, xshift=-2pt, yshift=-3pt] {$g_3$};

  % the chart atoms v_c = g_c/2, marked at half length on each arrow
  \foreach \px/\py in {0.3889/-0.4449, -0.6736/-0.5536,
                       0.6736/0.1773, -0.3889/0.8211}
    {\draw[line width=0.4pt, black!58, fill=white] (\px,\py) circle (1.4pt);}
  \node[font=\footnotesize, anchor=north east, inner sep=1.5pt]
    at (0.36,-0.49) {$v_0$};
  \fill (O) circle (1.1pt);
\end{scope}

\node[anchor=north west, align=left, font=\scriptsize, inner sep=0pt]
  at (-0.10,1.30)
  {$g_0=(1,1,1)$, \ $g_1=(1,-1,-1)$\\[3pt]
   $g_2=(-1,1,-1)$, \ $g_3=(-1,-1,1)$\\[3pt]
   $t_c=\tfrac14$, \ $\ell_c=\lVert g_c\rVert^{2}=3$\\[3pt]
   $s_c=t_c\ell_c=\tfrac34$, \ $\sum_cs_c=3=k$\\[3pt]
   $v_c=\sqrt{t_c}\,g_c=\tfrac12g_c$, \ $\lVert v_c\rVert^{2}=s_c$\\[3pt]
   $C=\{0,1,2\}$: the shaded face};

% ============ panel B: the two products of M_C, on R^3 ==================
\node[anchor=south, font=\small] at (9.20,4.90)
  {the two products of $M_C$: $k\times k$ and $\lvert C\rvert\times\lvert C\rvert$};

\node[font=\scriptsize, inner sep=1pt] at (6.00,3.60)
  {$M_C=\begin{pmatrix}1&1&1\\1&-1&-1\\-1&1&-1\end{pmatrix}$};
\node[anchor=north, align=center, font=\scriptsize, inner sep=3pt] at (6.00,2.95)
  {rows $g_c\tp$, $c\in C$\\ $\det M_C=4$};

\draw[ar] (7.80,3.87) -- (8.60,4.19);
\node[anchor=south east, font=\scriptsize, inner sep=1.5pt] at (8.62,4.21)
  {$M_C\tp M_C^{\vphantom{\mathsf{T}}}$};
\draw[ar] (7.80,3.33) -- (8.60,3.05);
\node[anchor=north east, font=\scriptsize, inner sep=1.5pt] at (8.62,3.00)
  {$M_C^{\vphantom{\mathsf{T}}}M_C\tp$};

\node[anchor=west, font=\scriptsize, inner sep=1pt] at (8.75,4.25)
  {$\SC=\begin{pmatrix}3&-1&1\\-1&3&1\\1&1&3\end{pmatrix}$};
\node[anchor=west, font=\scriptsize, inner sep=1pt] at (11.80,4.25)
  {$=4I_3-g_3^{\vphantom{\mathsf{T}}}g_3\tp$};
\node[anchor=west, font=\scriptsize, inner sep=1pt] at (8.75,2.95)
  {$\Gram_C=\begin{pmatrix}3&-1&-1\\-1&3&-1\\-1&-1&3\end{pmatrix}$};
\node[anchor=west, font=\scriptsize, inner sep=1pt] at (12.60,2.95)
  {$=4I_3-J_3$};

\node[anchor=north west, align=left, font=\scriptsize, inner sep=0pt]
  at (4.60,2.08)
  {$\tr\SC=\sum_{c\in C}\ell_c=9=\tr\Gram_C$\\[3.5pt]
   $\det\SC=\det\Gram_C=16=(\det M_C)^{2}$\\[3.5pt]
   $\operatorname{spec}\SC=\operatorname{spec}\Gram_C=\{1,4,4\}$\\[3.5pt]
   $\SC\psd I_3\iff\Gram_C\psd I_3\iff W[C]\psd0$\\[3.5pt]
   $\lmin(\SC)=1$: $C$ dominates, no slack};

% ================= the size strip at rank three ==========================
\begin{scope}[shift={(11.20,1.80)}]
  \node[anchor=east, font=\scriptsize, inner sep=3pt] at (-0.02,0.22)
    {$k=3$};
  \foreach \mm [count=\ii from 0] in {3,4,5,6,7,8,9,10}{%
    \draw[cellline] ({0.44*\ii},0) rectangle ++(0.44,0.44);
    \node[font=\tiny, anchor=north, inner sep=2.5pt] at ({0.22+0.44*\ii},-0.02)
      {\mm};}
  % the design sits at m = 4, the second cell
  \node[font=\tiny, inner sep=0pt] at (0.66,0.22) {$\dsn$};
  % the two cells left open at rank three, m = 6 and m = 7
  \draw[border, line width=0.9pt] (1.32,0) rectangle ++(0.88,0.44);
  % the window
  \draw[line width=0.7pt] (2.20,-0.34) -- (2.20,0.44);
\end{scope}

\node[anchor=north west, align=left, font=\scriptsize, inner sep=0pt]
  at (10.60,1.50)
  {$\win(k)=\tfrac{k(k+1)}{2}+1$, so $\win(3)=7$:\\[3pt]
   the statement at every $m\le\win(k)$\\[3pt]
   carries every $m$ (Theorem~\ref{thm:window}).};

% ============ panel C: the chart, on the index space R^4 ================
% The row heading stands at the left of the row and not above it: the panel-B
% list hangs to y = -0.35 across x 4.60..10.35 and the index rows of P and W
% sit at y = -0.46, so nothing of this height fits over the grids.  It is
% dropped below the box of the panel-A list, which ends at y = -1.69, and held
% above the box of the V caption, which begins at y = -2.44.
\node[anchor=east, font=\small] at (4.70,-2.12)
  {the index space $\R^{4}$};

\node[anchor=east, font=\scriptsize, inner sep=1pt] at (4.86,-1.48)
  {$\tfrac12$};
\cellframe{4.90}{-0.60}{3}{4}
\cellnums{4.90}{-0.60}{{1,1,1},{1,-1,-1},{-1,1,-1},{-1,-1,1}}
\node[anchor=north, align=center, font=\scriptsize, inner sep=3pt]
  at (5.56,-2.44) {$V$, rows $v_c\tp$\\ $V\tp V=I_3$};

\draw[ar] (6.42,-1.48) -- (8.02,-1.48);
\node[anchor=south, font=\scriptsize, inner sep=2pt] at (7.22,-1.42)
  {$P=VV\tp$};

\node[anchor=east, font=\scriptsize, inner sep=1pt] at (8.22,-1.48)
  {$\tfrac14$};
\cellframe{8.30}{-0.60}{4}{4}
\cellindex{8.30}{-0.60}
\cellnums{8.30}{-0.60}{{3,-1,-1,-1},{-1,3,-1,-1},{-1,-1,3,-1},{-1,-1,-1,3}}
\node[anchor=north, align=center, font=\scriptsize, inner sep=3pt]
  at (9.18,-2.44) {$P_{cd}=\sqrt{t_ct_d}\,\langle g_c,g_d\rangle$\\
                   $P_{cc}=s_c$, \ $\tr P=3=k$};

\draw[ar] (10.26,-1.48) -- (11.86,-1.48);
\node[anchor=south, font=\scriptsize, inner sep=2pt] at (11.06,-1.42)
  {$-\,D$};
\node[anchor=north, font=\scriptsize, inner sep=2pt] at (11.06,-1.54)
  {$D=\tfrac14I_4$};

\node[anchor=east, font=\scriptsize, inner sep=1pt] at (12.06,-1.48)
  {$\tfrac14$};
\path[fill=black!25] (12.14,-0.60) rectangle ++(1.32,-1.32);
\cellframe{12.14}{-0.60}{4}{4}
\cellindex{12.14}{-0.60}
\cellnums{12.14}{-0.60}{{2,-1,-1,-1},{-1,2,-1,-1},{-1,-1,2,-1},{-1,-1,-1,2}}
\node[anchor=north, align=center, font=\scriptsize, inner sep=3pt]
  at (13.02,-2.44) {$W=P-D$\\ $W_{cc}=t_c(\ell_c-1)$};

\node[anchor=north west, align=left, font=\scriptsize, inner sep=0pt]
  at (0.30,-3.43)
  {$W[C]=\tfrac14\bigl(3I_3-J_3\bigr)$, \quad
     $\operatorname{spec}W[C]=\{0,\tfrac34,\tfrac34\}$, \quad
     $D=\diag(t)$,
     $\sum_cW_{cc}=k-1$, \quad $\pi_C=\det P[C]=\tfrac{16}{64}=\tfrac14$\\[3.5pt]
   $\det W[C]=\bigl(\prod_{c\in C}t_c\bigr)\det\bigl(\Gram_C-I_3\bigr)
     =\tfrac1{64}\cdot0=0$ \ by Corollary~\ref{cor:minor}, \ while
     $\det W[\{0,1\}]=\tfrac1{16}\cdot3=\tfrac3{16}$};

\end{tikzpicture}
